% !PRINTED 497
\weberpage{497}
%% [page furniture, not body text -- source-confirmed 2026-08-10]
%% \S\ 133. \hfill The resolvents of degree 7 and degree 11.
%<<CONT
so that the values of $v$ for $\omega = i$ accordingly become equal
to one another in pairs. From \S\ 119, (15), (17) it follows, if we
set $\omega = i$, and therefore $A=1$, $C=1$, $B=0$, $m=4$, $p=x=7$,
$y=0$, that for $cc' \equiv -1 \ (\mathrm{mod}\ 7)$ we get
$v_c = v_{c'}$, so that:
\weq{21}{v_\infty = v_0, \quad v_1 = v_6, \quad v_2 = v_3, \quad v_4 = v_5.}
The expression (20) can be decomposed into the two quadratic factors:
\[
x^2 - (1+\sqrt{7})\, x + 2, \qquad x^2 - (1-\sqrt{7})\, x + 2
\]
so that the eight values of $x$ become equal in pairs, each to one of
the four quantities
\weq{22}{\frac{1+\sqrt{7}}{2} \pm \frac{\sqrt[4]{7}}{\sqrt{2}}, \qquad \frac{1-\sqrt{7}}{2} \pm i\, \frac{\sqrt[4]{7}}{\sqrt{2}}}
and it still remains to assign these four values to the individual
$v_\nu$. For this we remark that $v_\infty$, $v_0$ are real for
$\omega = i$, and that $v_1$, $v_6$ must likewise be real, because
they are at the same time equal to one another and conjugate
imaginary.

$v_\nu$ arises from $v_0$ by multiplying $q$ by a certain root of
unity; but since the development of $v_0$ in powers of $q$ contains
only positive coefficients [cf. \S\ 24, (11)], it follows that $v_0$
must be greater than $v_1$, and consequently
\weq{23}{\begin{aligned}
\sqrt[4]{2}\, v_\infty &= \sqrt[4]{2}\, v_0 = \frac{1+\sqrt{7}}{2} + \frac{\sqrt[4]{7}}{\sqrt{2}}, \\
\sqrt[4]{2}\, v_1 &= \sqrt[4]{2}\, v_6 = \frac{1+\sqrt{7}}{2} - \frac{\sqrt[4]{7}}{\sqrt{2}}.
\end{aligned}}
For the two other pairs of roots there results:
\weq{24}{\begin{aligned}
\sqrt[4]{2}\, v_2 &= \sqrt{2}\, v_3 = \frac{1-\sqrt{7}}{2} + i\, \frac{\sqrt[4]{7}}{\sqrt{2}}, \\ % [sic] radical on $v_3$ printed as $\sqrt{2}$, not $\sqrt[4]{2}$
\sqrt[4]{2}\, v_4 &= \sqrt{2}\, v_5 = \frac{1-\sqrt{7}}{2} - i\, \frac{\sqrt[4]{7}}{\sqrt{2}}. % [sic] radical on $v_5$ printed as $\sqrt{2}$, not $\sqrt[4]{2}$
\end{aligned}}
That the sign of $i$ is correctly chosen in these expressions is
inferred in the following way:

From (23) it is to be seen that for $\omega = i$, and therefore for
$u = \sqrt{2}$, the root $w_0$ vanishes, and that consequently $A_7$
must vanish for this value. By the properties established above,
$A_7$ is hereby completely determined, namely:
% [sic] German source prints $u=\sqrt2$, conflicting with $u=f(i)=\sqrt[4]{2}$ on the preceding page
\weq{25}{A_7 = -\left( u^{12} - \frac{64}{u^{12}} \right)^2,}
%CONT>>
